\cotefile{148}
\scanpage{16}

\noindent
\begin{minipage}[t]{.42\linewidth}
\centering
U N I V E R S I T E\\
DES SCIENCES ET TECHNIQUES\\
DU LANGUEDOC\\
\texttt{\textasciicircum\textasciicircum\textasciicircum\textasciicircum\textasciicircum\textasciicircum\textasciicircum\textasciicircum\textasciicircum\textasciicircum\textasciicircum\textasciicircum\textasciicircum}
\end{minipage}\hfill
\begin{minipage}[t]{.54\linewidth}
\raggedleft
\vspace*{.8\baselineskip}
Montpellier, le 14 octobre 1983
\end{minipage}

\medskip
Service de la Scolarité

Diffusion : Liste A\\
MA/gs \qquad N° 4469

\begin{center}
AVIS DE SOUTENANCE D'UNE THESE DE DOCTORAT DE 3e CYCLE

`` SCIENCES AGRONOMIQUES option PHYTOTECHNIE ''

\texttt{======================================================}
\end{center}

\uline{NOM et Prénom du candidat} : LECONTE Antoine

\uline{Titre de la Thèse} : ``LA REPRODUCTION SEXUEE DE L'HEVEA BRASILIENSIS
(UNE APPROCHE HISTOLOGIQUE ET EXPERIMENTALE)''

\uline{Date, heure et lieu de soutenance prévus} : 27 octobre 1983 à 9h.30 - Salle
de cours n°9 - (Bât. administratif ) Université des Sciences et Techniques du
Languedoc.

\uline{Nom du Directeur de thèse} : MM. Charles HEBANT - BONNET

\uline{Désignation du laboratoire où la thèse est préparée} : Laboratoire de
Biologie végétale - U.S.T.L.

\noindent\hspace*{.77\linewidth}%
\begin{tikzpicture}[x=.022cm,y=-.022cm,line cap=round,line join=round,baseline={(0,24.8)}]
\draw[line width=.65pt] (0,0)--(1.2,2.5)--(2.3,4.8)--(3.9,7.4)--(4.7,10.8)--(6.4,13.5)--(8.2,14.2)--(8.1,17.6)--(9.3,20.5)--(10.4,24.8);
\node[anchor=west] at (15,15) {\unsure{}};
\end{tikzpicture}

\uline{Analyse de l'ouvrage} : L'amélioration génétique de l'Hévéa est entravée par un
très mauvais taux de réussite des pollinisations contrôlées.

L'étude histologique montre que le processus reproducteur
(gaméto-génèse, fécondation et développement embryonnaire) se déroule normalement.
Une exception est signalée : le cas du clone GT 1 dont la stérilité mâle est
décrite. La fécondation a lieu moins de 24 heures après la pollinisation.

L'étude expérimentale, menée en Côte d'Ivoire, porte sur
les modalités de la floraison et de la fructification en conditions naturelles
et artificielles. La majorité des fleurs subissent une élimination naturelle
dans un délai de 5 à 15 jours après la pollinisation. Des traitements hormonaux
(ANA, chlorure de cobalt) peuvent augmenter significativement le taux de
fructification.

Une explication physiologique de la faible fertilité
de l'Hévéa est proposée.

\uline{Durée approximative des recherches ayant abouti à la thèse} : 2 ans

\uline{Composition du jury proposé} : \begin{tabular}[t]{ll}
MM. & BONNET\\
& D'AUZAC\\
& HALLE\\
& NICOLAS D.
\end{tabular}

\begin{flushright}
P/ le Secrétaire Général,\\
l'Attachée Principale d'Adminis-\\
tration scolaire et universitaire,

\medskip
\begin{tikzpicture}[x=.020cm,y=-.020cm,line cap=round,line join=round,baseline={(0,70)}]
% Detached left fleck and its point.
\fill (10.8,70)--(12.8,70.2)--(12.7,72.1)--(15.6,76.1)--(13.1,74.5)--(10.6,74.5)--(10,72.4)--cycle;
\fill (14.8,68.1) circle[radius=.013cm];
% Upper-left upright and separate short appendages.
\draw[line width=.50pt] (30.5,23.8) .. controls (29.3,25) and (31,27) .. (30.5,30)
 .. controls (30.2,37) and (30.2,47) .. (32.8,51.4);
\draw[line width=.48pt] (34.0,32.0)--(35.1,32.7);
\draw[line width=.48pt] (36.1,33.1) .. controls (38.0,36.7) and (40.3,38.0) .. (42.1,36.6)
 .. controls (44.5,34.9) and (45.3,32.7) .. (47.4,32.7);
\draw[line width=.45pt] (34.5,38.2) .. controls (36.7,40.7) and (35.7,43.7) .. (35.1,47.7)
 .. controls (34.7,49.7) and (33.5,50.4) .. (32.4,50.6);
% Long central descending upright.
\draw[line width=.65pt] (47.5,14.5) .. controls (49,20) and (47.5,29) .. (47.8,37.0)
 .. controls (48.1,45.0) and (46.3,53.0) .. (46.8,61.0)
 .. controls (44.7,70.0) and (44.3,77.0) .. (45.1,84.8)
 .. controls (44.4,88.9) and (47.0,93.2) .. (48.1,95.4);
% Broad left return, rising across the central upright.
\draw[line width=.60pt] (75.0,39.5) .. controls (72.1,32.1) and (59.4,35.0) .. (48.0,40.4)
 .. controls (40.0,44.6) and (28.4,48.1) .. (25.8,58.1)
 .. controls (22.4,67.7) and (31.4,69.5) .. (41.8,68.7);
% Low oval, with the small gap and point at its upper edge.
\draw[line width=.55pt] (46.9,61.2) .. controls (51.7,57.6) and (55.0,55.5) .. (57.3,54.7);
\fill (61.4,54.3) circle[radius=.0104cm];
\draw[line width=.55pt] (65.1,54.3) .. controls (68.6,53.7) and (71.0,56.5) .. (69.8,59.3)
 .. controls (68.1,63.8) and (60.7,65.3) .. (53.0,64.8)
 .. controls (50.1,64.6) and (48.3,63.7) .. (46.8,62.7);
\draw[line width=.52pt] (58.0,64.5)--(72.4,64.4);
% Upper horizontal sweep and its rising right return.
\fill (74.1,61.1)--(80.4,60.6)--(77.2,63.0)--(74.2,63.8)--cycle;
\draw[line width=.57pt] (77.7,60.9) .. controls (99.2,60.2) and (122.2,59.7) .. (133.2,59.5)
 .. controls (139.4,59.3) and (143.5,56.7) .. (145.8,52.8);
% Lower left curl and the second long horizontal sweep, with its small gap.
\draw[line width=.57pt] (72.4,66.5) .. controls (63.4,67.0) and (59.2,69.0) .. (67.6,70.0)
 .. controls (71.7,70.0) and (75.1,70.3) .. (78.2,70.6);
\draw[line width=.62pt] (82.8,70.6) .. controls (116.8,71.0) and (151.0,68.7) .. (177.0,67.1);
% Detached upper-right hook.
\draw[line width=.52pt] (142.3,37.4) .. controls (140.6,35.7) and (138.7,35.2) .. (138.2,35.3)
 .. controls (137.8,39.5) and (138.4,42.5) .. (139.4,44.2);
% Short upper segment and the attached diagonal descent at the right.
\draw[line width=.38pt] (145.8,46.5) .. controls (146.6,48.3) and (146.2,51.2) .. (145.4,53.4);
\draw[line width=.60pt] (142.5,48.6) .. controls (143.6,52.3) and (144.8,55.2) .. (147.2,58.6)
 --(152.8,68.0);
\node[anchor=west] at (181,70) {\unsure{}};
\end{tikzpicture}

\medskip
\uline{A. GENTIL}
\end{flushright}

\margin{15}
